\section{A local criterion for nontrivial loops}\label{sec:paths}

The geometric construction will attach group labels to paths only within
bounded regions.  We therefore need a criterion that detects a nontrivial
loop using such local information.  The criterion below uses hyperbolic
graphs as the local models.  It does not require a global map to a model,
local finiteness of a graph, or a bound on the information carried by one
edge.

A model to keep in mind is a graph whose vertices are points of a space
and whose edges are paths confined to small regions.  A reading lifts
such paths to a covering space chosen near the initial vertex; its
endpoint records both the actual endpoint and the sheet of the cover.
The argument follows the coarse local-to-global viewpoint of hyperbolic
geometry \cite{Gromov1987}; we prove here the version for bounded path
readings in compatible charts.

\subsection{Paths and their readings}

All graphs have oriented edges with specified reverses; multiple edges and
loops are allowed.  Each edge has length one.  Paths are finite edge paths,
and their lengths count edges.  A constant path has length zero.  The path
groupoid of a graph is obtained by cancelling consecutive reverse edges.
We use the path metric on each connected component of a graph.

Fix a positive integer $N$.  At each vertex $v$ of a graph $G$, suppose that
we are given a pointed graph $(D_v,d_v)$ and a rule for reading every path
of length at most $N$ starting at $v$ as an edge path in $D_v$ starting at
$d_v$.  Write $[p]_v$ for the terminal vertex of the reading of $p$.
The following are the required properties.
\begin{enumerate}
\item[\textnormal{(P1)}] Readings respect prefixes and reverses of edges,
and read a backtrack as a backtrack.  If $p$ starts at $v$ and $|p|<N$,
reading gives a bijection between the edges leaving the actual endpoint of
$p$ and the edges leaving $[p]_v$.  If $[p]_v=[p']_v$, then $p$ and $p'$
have the same actual endpoint in $G$, and their edge bijections agree
whenever both are defined.
\item[\textnormal{(P2)}] Suppose that $p$ runs from $v$ to $w$ and that
$a,b$ start at $w$.  If each of $p,a,b$ has length at most $N/2$, then
\begin{equation}\label{path:change-root}
 [pa]_v=[pb]_v\quad\Longleftrightarrow\quad [a]_w=[b]_w.
\end{equation}
\item[\textnormal{(P3)}] All connected components of all the graphs $D_v$
are Gromov hyperbolic with a common constant.
\end{enumerate}
These rules concern paths within the stated range.  They do not assert
that the whole graph $G$ maps to $D_v$, or that the whole graph $D_v$ maps
to $G$.

A loop $p$ at $v$ \emph{closes on reading} if $[p]_v=d_v$.  Let
$\mathcal Q$ be the quotient of the path groupoid of $G$ by the relations
that every such loop of length at most three is the constant path.

\begin{lemma}[Local path criterion]\label{thm:path-lemma}
There is a bound $N_0$, depending only on the common hyperbolicity
constant, with the following property.  If the readings satisfy
\textnormal{(P1)--(P3)} with $N\ge N_0$, and $e$ is a single-edge loop at
$v$ satisfying $[e]_v\ne d_v$, then $e$ is nontrivial in $\mathcal Q$.
\end{lemma}

The proof constructs a locally geodesic representative for each path,
unique up to a uniformly narrow comparison.  Such representatives can be
continued one edge at a time.  The comparison class is unchanged by the
relations defining $\mathcal Q$, but distinguishes the edge in the lemma
from the constant path.

\subsection{Calculations within one chart}

We first record precisely how the reading rules allow short calculations
to move between charts.  By (P1), any edge path in $D_v$ starting at a
state already reached by reading can be lifted to a path in $G$, as long
as the total reading length is at most $N$.  If two such readings have the
same endpoint, their lifted paths have the same actual endpoint, and any
further common continuation reads identically.  Together with (P2), this
has three useful consequences.

First, a short circuit closes on reading if and only if its two
complementary paths from one corner to another have equal read endpoints.
Indeed one may read one path and the reverse of the other, using (P1).
Closure can be tested at any corner: a cyclic change of starting corner
amounts to prefixing by a side, applying (P2), and cancelling that side
using (P1).  Here and below, ``short'' means that all prefixes and
continuations in the comparison lie within the ranges of (P1) and (P2).

Second, geodesicity of a short reading is independent of a short prefix.
For example, if the reading of $a$ from $w$ is geodesic but its reading
after $p:v\to w$ has a shorter competitor, lift that competitor from
$[p]_v$.  It gives a path $b$ from $w$ with
$[pa]_v=[pb]_v$.  Equation~\eqref{path:change-root} gives
$[a]_w=[b]_w$, contradicting geodesicity.  The reverse implication uses a
shorter competitor in $D_w$ and the same argument.  For distance
comparison, suppose the states are reached by paths $a,b$ from $w$.
Lift a shortest connector $c$ between their read states, starting after
$a$ (or after $pa$).  Its length is at most $|a|+|b|$, and apply (P2)
to $ac$ and $b$.  Doing this in both charts proves equality of the
distances whenever $|p|\le N/2$ and $2|a|+|b|\le N/2$.
These bounds, as well as their versions with $a,b$ interchanged, hold
in every use below.

Third, a short strip of closed circuits can be calculated in one chart
without adding the lengths of all its transverse paths.  Here is the
explicit bookkeeping.  Let $a_j,b_j$ be the prefixes down the two sides
of the strip, and let $\rho_j$ be its transverse path from the endpoint
of $a_j$ to that of $b_j$.  Starting in the chart at the beginning of the
first side, the equalities to prove are
\begin{equation}\label{path:strip-calculation}
 [a_j\rho_j]=[\rho_0b_j].
\end{equation}
Closure of the next elementary circuit, transported after $a_j$ by (P2),
compares its two routes across the next step.  Explicitly, write
$a_{j+1}=a_js_j$ and $b_{j+1}=b_jt_j$, where $s_j,t_j$ are the next
side portions.  Closure and then the preceding equality give
\[
 \begin{aligned}
 [a_{j+1}\rho_{j+1}]&=[a_j\rho_jt_j]\\
                    &=[\rho_0b_jt_j]=[\rho_0b_{j+1}],
 \end{aligned}
\]
using (P1) for the common continuation, edge by edge.  This proves
\eqref{path:strip-calculation} at the next step.  Thus only the side
prefixes and each individual transverse path enter the range bound.
Conversely, transverse paths found in this one chart lift to $G$ and
give circuits that close when read from their own corners.  These facts
will justify all uses of diagrams below.

Choose an integer $\Delta\ge10$ large enough for the following standard
consequences of the common hyperbolicity bound.  In every component of
every chart, geodesic triangles are $\Delta$-slim on vertices: every
vertex on one side lies within $\Delta$ of a vertex on one of the other
two sides.  Dividing a geodesic quadrangle by a diagonal then shows that
each side lies within $2\Delta$ of the union of the other three sides.
The Gromov product
\[
 (x\mid y)_u=\tfrac12\bigl(d(u,x)+d(u,y)-d(x,y)\bigr)
\]
satisfies
\begin{equation}\label{path:product}
 (x\mid z)_u\ge
 \min\{(x\mid y)_u,(y\mid z)_u\}-\Delta.
\end{equation}
Moreover, for a geodesic segment $[x,z]$, there is a vertex on that segment
at distance at most $(x\mid z)_u+10\Delta$ from $u$.  Enlarging $\Delta$
absorbs all rounding to vertices.  These are the usual equivalent forms
of hyperbolicity; see \cite[Chapter III.H]{BridsonHaefliger}.

We will also use the following elementary comparison.  If two geodesic
segments $\gamma_1,\gamma_2$, of lengths $\ell_1,\ell_2$, have
corresponding endpoints at distance at most $h$, then for every integer
$t\ge0$,
\begin{equation}\label{path:geodesic-comparison}
 d\bigl(\gamma_1(\min\{t,\ell_1\}),
        \gamma_2(\min\{t,\ell_2\})\bigr)\le10(h+\Delta).
\end{equation}
This is a \emph{synchronous comparison} of width $10(h+\Delta)$:
the parameters advance by one on both segments until the shorter
segment ends, and then that segment waits at its endpoint.
To check the bound, join the
corresponding endpoints by geodesics of length at most $h$.  The resulting
quadrangle is $2\Delta$-slim.  A point close to the opposite geodesic has
a mate whose distance from that geodesic's initial endpoint differs from
its own parameter by at most $h+2\Delta$.  A point close to an end
connector is within $h+2\Delta$ of the corresponding end.  These estimates,
together with the bound $2h$ on the difference of the lengths, imply
\eqref{path:geodesic-comparison}, also when one segment has ended.

Fix integers
\begin{equation}\label{path:constants}
 H=1000\Delta,\qquad L\ge10^6H,\qquad k=20L,
 \qquad N\ge10^5k.
\end{equation}
Taking $L=10^6H$, for example, gives a threshold $N_0=10^5(20L)$
depending only on the original hyperbolicity constant.
A \emph{guide} is a path each of whose subpaths of length at most $k$
reads as a geodesic in the chart at its initial vertex.  Its entire
length may be arbitrarily large.  Every individual chart calculation
below uses paths, including prefixes, of length less than $1000k$.
Consequently (P1) and (P2) always apply.  A long guide is handled in
separate short pieces, never by reading its whole prefix.

\subsection{Narrowing comparisons between guides}

For paths $\alpha,\beta$ with the same initial and terminal vertices, a
\emph{ladder of width $h$} consists of the following data.  A finite
schedule of pairs of positions begins at $(0,0)$ and ends at
$(|\alpha|,|\beta|)$; at each step each position advances by zero or one.
At each scheduled pair there is a path, called a rung, from the vertex
on $\alpha$ to the vertex on $\beta$, of length at most $h$.  The first
and last rungs are empty.  Each elementary circuit, formed by consecutive
rungs and the intervening portions of the two paths, must close on
reading.  Pausing both paths is allowed.  In particular, this definition
does not require an elementary circuit to have length at most three.
Ladders are comparisons of readings, rather than expressions in the
relations defining $\mathcal Q$.

For the widths used below, ladders are symmetric, by changing corners of
their elementary circuits.  They compose with addition of widths whenever
the sum is at most $1000H$, which suffices for every composition below.
Indeed, given ladders from
$\alpha$ to $\beta$ and from $\beta$ to $\gamma$, synchronize the advances
along $\beta$, inserting pauses in the other sides as necessary, and
concatenate the two rungs.  Each resulting circuit is a union of one or
two elementary circuits from the given ladders, so closes by
\eqref{path:strip-calculation}.  Appending the same trailing path to both
sides preserves the width: use empty rungs along that trailing path.

\begin{lemma}[Narrowing]\label{path:narrowing}
If two guides have a ladder of width at most $1000H$, then they have one
of width at most $H$.
\end{lemma}

\begin{proof}
Let $\alpha,\beta$ be the guides, and let the given width be
$h\le1000H$.  Consider an interval of the schedule on which $\alpha$
advances by at most $8L$.  On this interval $\beta$ advances by less than
$k$.  Otherwise stop at its first advance of $k$ edges.  The strip up to
that point calculates in one chart: its sides have lengths at most
$8L$ and $k$, and its rungs have length at most $h$.  The second side is
geodesic there, whereas its endpoints can be joined using the first side
and the two end rungs.  This gives the contradiction
\[
 k\le8L+2h<20L=k.
\]
It follows also that on every such schedule interval both sides read as
geodesics, and their lengths differ by at most $2h$.

Suppose first that $|\alpha|\ge L$.  Mark $\alpha$ at
$x_0,\ldots,x_m$, including its endpoints, so that consecutive marked
intervals have lengths in $[L,2L]$.  Choose scheduled mates
$b_0,\ldots,b_m$ on $\beta$, using its actual endpoints for $b_0,b_m$.
The distance along $\beta$ between consecutive old mates is at least
$L-2h$.

\begin{figure}[htbp]
\centering
\begingroup
\definecolor{ladderaccent}{RGB}{22,91,119}
\begin{tikzpicture}[x=1cm,y=1cm,>=Latex,
  every node/.style={font=\small},
  guide/.style={line width=.9pt},
  oldrung/.style={densely dashed,line width=.65pt,black!65},
  newrung/.style={line width=1.1pt,ladderaccent}]
  \coordinate (am) at (0,2.05);
  \coordinate (ai) at (5.4,2.05);
  \coordinate (ap) at (10.8,2.05);
  \coordinate (bm) at (0,.2);
  \coordinate (bo) at (6.7,.2);
  \coordinate (bn) at (5.4,.2);
  \coordinate (bp) at (10.8,.2);
  \draw[guide] (am)--(ap);
  \draw[guide] (bm)--(bp);
  \draw[oldrung] (am)--node[left] {$\le h$} (bm);
  \draw[oldrung] (ap)--node[right] {$\le h$} (bp);
  \draw[oldrung] (ai)--node[pos=.50,right,fill=white,inner sep=1.5pt]
    {old: $\le h$} (bo);
  \draw[newrung] (ai)--node[pos=.55,left,fill=white,inner sep=1.5pt]
    {new: $\le 2\Delta$} (bn);
  \foreach \p in {am,ai,ap,bm,bo,bp}
    \fill (\p) circle (1.7pt);
  \draw[newrung,fill=white] (bn) circle (2.2pt);
  \node[above=4pt] at (am) {$x_{i-1}$};
  \node[above=4pt] at (ai) {$x_i$};
  \node[above=4pt] at (ap) {$x_{i+1}$};
  \node[below=5pt] at (bm) {$b_{i-1}$};
  \node[below=5pt] at (bn) {$b_i'$};
  \node[below=5pt] at (bo) {$b_i$};
  \node[below=5pt] at (bp) {$b_{i+1}$};
  \node[left=18pt] at (am) {$\alpha$};
  \node[left=18pt] at (bm) {$\beta$};
  \draw[decorate,decoration={brace,amplitude=4pt}]
    (0,2.76)--node[above=5pt] {length in $[L,2L]$} (5.4,2.76);
  \draw[decorate,decoration={brace,amplitude=4pt}]
    (5.4,2.76)--node[above=5pt] {length in $[L,2L]$} (10.8,2.76);
\end{tikzpicture}
\endgroup
\caption{Narrowing a ladder in a single chart. Both displayed guide
  pieces read as geodesics. The interior mark $x_i$ is far from the
  end rungs, so quadrangle slimness supplies a new mate $b_i'$ on
  the opposite side, within $2\Delta$ of $x_i$. The old scheduled
  mate $b_i$ is at distance at most $h+2\Delta$ from $b_i'$ along
  that side. The drawing is
  schematic: the long side intervals and the two rung bounds use
  different visual scales.}
\label{fig:path-ladder}
\end{figure}


At an interior mark $x_i$, calculate the part between $x_{i-1}$ and
$x_{i+1}$ in one chart, as in Figure~\ref{fig:path-ladder}.
Replace its end rungs by geodesics of length at
most $h$.  The point $x_i$ is at distance at least $L-h>2\Delta$ from
either end connector.  Quadrangle slimness therefore puts it within
$2\Delta$ of a vertex $b'_i$ on the opposite geodesic, the corresponding
piece of $\beta$.  Its old mate $b_i$ was within $h$ of $x_i$, so
geodesicity gives
\begin{equation}\label{path:mate-shift}
 |b'_i-b_i|_{\beta}\le h+2\Delta.
\end{equation}
Here the left side denotes distance in the parameter of the path
$\beta$.  The new mates remain in order, since the difference between
consecutive new positions is at least
\[
 L-2h-2(h+2\Delta)=L-4h-4\Delta>0.
\]
The same estimate at the ends puts them strictly between the unchanged
endpoint mates.  Lift short paths from $x_i$ to $b'_i$ to obtain new
rungs of length at most $2\Delta$; use empty rungs at $i=0,m$.

The new rungs agree with the old calculation: following an old rung and
the relevant portion of $\beta$ has the same read endpoint as the new
rung.  On a marked interval this equality can be checked together with
both adjacent marked intervals.  Their union has first-side length at
most $6L$, so the preceding bound and the short-strip calculation apply.
Thus between successive new mates both sides are geodesic, and their
end rungs have length at most $2\Delta$.  Apply
\eqref{path:geodesic-comparison}, lift the comparison rungs, and retain
the specified rungs at the marked endpoints.  Each elementary circuit
closes in this chart and hence at its own corner.  The resulting width
is at most $30\Delta<H$.

If $|\alpha|<L$, the initial schedule argument applies to the whole
ladder.  Both sides calculate as geodesics in one chart with identical
endpoints, and \eqref{path:geodesic-comparison} directly gives width at
most $10\Delta<H$.
\end{proof}

In particular, existence of a width-$H$ ladder is an equivalence relation
on guides with fixed endpoints.  Reflexivity uses empty rungs, symmetry
was noted above, and transitivity follows by composing two ladders and
applying Lemma~\ref{path:narrowing}.

\subsection{Continuing a guide by one edge}

We next show that this comparison class can be continued along any edge.
The difficulty is that appending an edge need not preserve local
geodesicity.  We repair the guide by choosing nearby vertices at widely
spaced marks and minimizing the total distance between those choices.
All nearby choices are specified by paths in the charts; distance in
$G$ alone would forget the states that must be compared.

\begin{lemma}[Extension]\label{path:extension}
If $\alpha$ is a guide and $e$ is an edge starting at its terminal vertex,
there is a guide $\beta$ to the endpoint of $\alpha e$ and a ladder
between $\alpha e$ and $\beta$ of width at most $50H$.
\end{lemma}

\begin{proof}
If $|\alpha|<L$, read $\alpha e$ in its initial chart and join its
endpoints by a geodesic.  Lift that geodesic to $G$.  The result is a
guide by transport of geodesicity.  Compare it with $\alpha$, whose
final endpoint differs by at most one in the chart, using
\eqref{path:geodesic-comparison}.  Finish the comparison by letting the
first side traverse $e$ while the second waits.  This proves the claim
in this case.

Suppose $|\alpha|\ge L$, and mark it at
$a_0,\ldots,a_m$ in intervals $\alpha_i$ of lengths
$\ell_i\in[L,2L]$, where $\alpha_i$ runs from $a_i$ to $a_{i+1}$.
At each mark choose a path $c_i$ starting there, with $|c_i|\le H$.
Require $c_0$ to be empty and $c_m=e$.  For adjacent marks define
\[
 F_i(c_i,c_{i+1})=
 d_{D_{a_i}}\bigl([c_i]_{a_i},[\alpha_i c_{i+1}]_{a_i}\bigr).
\]
Choose the paths $c_i$ to minimize $\sum_{i=0}^{m-1}F_i$.  A minimum
exists because the set of choices is nonempty and the objective takes
nonnegative integer values; no compactness assumption is needed.

The point of this minimization is that, on each short block,
hyperbolicity will supply a geodesic in one chart between its chosen
outer endpoints, with points lying in order within $H$ of the interior
marks.  Replacing the interior choices by short paths to those points
will force equality in the triangle inequality for the minimizing
choices.

For each $i$, lift a geodesic realizing $F_i$, starting after $c_i$ in
the chart at $a_i$, to a path $\beta_i$ in $G$.  Equality of read
endpoints shows that $\beta_i$ ends at the actual endpoint of $c_{i+1}$.
Thus the paths concatenate to a path
$\beta=\beta_0\cdots\beta_{m-1}$ from $a_0$ to the endpoint of $e$.
The comparison circuit between $\alpha_i$ and $\beta_i$, using
$c_i,c_{i+1}$, closes on reading.  The paths $c_i$, rather than just
their actual endpoints, retain the states used in this assertion.

We claim that on every block of at most $100$ consecutive marked
intervals the concatenation of the $\beta_i$ is geodesic in a single
chart calculation.  Read such a block in the chart at its first mark.
The paths $\alpha_i$, $c_i$, and $\beta_i$ all calculate there
consistently, by \eqref{path:strip-calculation}; each side has length
bounded by $100(2L+2H)$, well within the prescribed range.  Distances
computed here agree with the single-interval distances $F_i$, by
transport of short competitors.

For the hyperbolic calculation, denote the read marks on this block by
$x_0,\ldots,x_q$, where $q\le100$.  Their consecutive distances are in
$[L,2L]$, and
\begin{equation}\label{path:zero-products}
 (x_{i-1}\mid x_{i+1})_{x_i}=0\qquad(1\le i<q),
\end{equation}
because two consecutive marked intervals of $\alpha$ have total length
at most $4L<k$.  We record the consequences of
\eqref{path:zero-products} explicitly.  Induction using
\eqref{path:product} gives
\begin{equation}\label{path:forward-product}
 (x_0\mid x_{i+1})_{x_i}\le\Delta\qquad(1\le i<q).
\end{equation}
The first case follows from \eqref{path:zero-products}.  For the next
case, the preceding estimate and the identity
\[
 (x_0\mid x_{i-1})_{x_i}
 =d(x_{i-1},x_i)-(x_0\mid x_i)_{x_{i-1}}
 \ge L-\Delta
\]
show that this product exceeds $\Delta$.  Applying
\eqref{path:product} to the zero product
$(x_{i-1}\mid x_{i+1})_{x_i}$ forces
\eqref{path:forward-product}.  The same argument backwards, followed
by another application of \eqref{path:product}, gives
\begin{equation}\label{path:block-products}
 (x_0\mid x_q)_{x_i}\le2\Delta\qquad(1\le i<q).
\end{equation}
For the last implication, use
$(x_{i-1}\mid x_q)_{x_i}\le\Delta$ and
$(x_{i-1}\mid x_0)_{x_i}\ge L-\Delta>2\Delta$.
Also, the distances from $x_0$ increase at every mark by at least
$L-2\Delta$, since
\[
 d(x_0,x_{i+1})-d(x_0,x_i)
 =d(x_i,x_{i+1})-2(x_0\mid x_{i+1})_{x_i}.
\]
The analogous estimate holds from the other endpoint.

Let $A,B$ be the read endpoints of the chosen paths at the two ends of
the block.  They are at distance at most $H$ from $x_0,x_q$,
respectively.  By \eqref{path:block-products}, every interior $x_i$
has a mate on a geodesic $[x_0,x_q]$ at distance at most $12\Delta$.
Such a mate is farther than $2\Delta$ from either of the geodesic
connectors $[x_0,A]$ and $[x_q,B]$: its distance from those connectors
is at least $L-H-14\Delta$.  Quadrangle slimness therefore gives a
vertex $y_i$ on $[A,B]$ with
\begin{equation}\label{path:block-mates}
 d(x_i,y_i)\le20\Delta<H.
\end{equation}
The vertices $y_i$ occur in increasing order along $[A,B]$.  Indeed
their distances from $A$ differ from the corresponding distances
$d(x_0,x_i)$ by at most $H+20\Delta$.  Consequently each consecutive
difference is at least
\[
 L-2\Delta-2H-40\Delta>0.
\]
The same inequalities keep them strictly between $A$ and $B$.

Lift short paths from the interior $x_i$ to the $y_i$ to replace the
interior choices $c_i$, keeping the two outer choices fixed.  These are
legal choices by \eqref{path:block-mates}.  For them the sum of the
successive distances equals $d(A,B)$, since their endpoints occur in
order on a geodesic.  Transport of distances shows that this is also
their sum of single-interval objectives $F_i$.  Minimality of the
original total objective implies that its sum on this block is at most
$d(A,B)$: changing only interior choices leaves every term outside the
block unchanged.  The triangle inequality gives the reverse inequality.
Thus the original joins concatenate geodesically on the block, as
claimed.

Each join $\beta_i$ has length at least $L-2H$.  A subpath of $\beta$
of length at most $k=20L$ therefore meets fewer than $100$ consecutive
joins: apart from at most two end joins, every join it meets is entirely
contained in that subpath.  The block result and transport of
geodesicity prove that $\beta$ is a guide.

Finally compare $\alpha_i$ and $\beta_i$ in their single-interval
chart.  Their corresponding endpoints are joined by the given paths
$c_i,c_{i+1}$ of length at most $H$.  Formula~\eqref{path:geodesic-comparison}
gives comparison rungs of length at most $10(H+\Delta)$; retain the
given choice paths as rungs at the marks.  Lift these comparisons and
concatenate them.  At the last mark the rung is $e$; let the first
side traverse $e$ and finish with an empty rung.  This gives the required
ladder of width at most $50H$.
\end{proof}

\subsection{The invariant of a path}

We can now finish the local criterion.  Fix an initial vertex $v$ and
consider guides starting at $v$, modulo width-$H$ ladders.  By
Lemma~\ref{path:extension}, appending an edge $e$ defines a continuation
of such a class.  It is independent of both choices involved.  If
$\alpha,\alpha'$ represent the same class, append $e$ to their
width-$H$ ladder.  If $\beta,\beta'$ are any two guides supplied by
Lemma~\ref{path:extension}, composition gives a ladder between them of
width at most
\[
 50H+H+50H=101H.
\]
Lemma~\ref{path:narrowing} reduces this to width $H$.  The same argument
covers different choices of the extension for a fixed guide.

Let $q$ be a loop of length at most three that closes on reading.
It has a ladder of width at most three with the constant path.  One
explicit schedule traverses $q$ on the first side while the second
waits: after a noninitial position use the remaining suffix of $q$ as
the rung, and use empty first and last rungs.  The first circuit closes
because $q$ does, and the subsequent ones cancel consecutive reverse
edges.  The short-chart rules justify the same assertion at each corner.
Appending this ladder after any guide $\alpha$ gives a width-three
ladder between $\alpha q$ and $\alpha$.

Successively applying Lemma~\ref{path:extension} along $q$ produces a
guide $\beta$ with a ladder to $\alpha q$ of width at most $150H$.
To see the bound, at each step append the remaining common trailing
edges to the comparison and compose; each of at most three steps adds
$50H$.  Composing with the preceding width-three ladder and narrowing
shows that $\beta$ and $\alpha$ represent the same class.  The same
reasoning applies to a backtrack, which closes on reading by (P1).

Starting from the empty guide at $v$ and continuing along a path thus
gives a class unchanged by every defining relation of $\mathcal Q$.
Insertion of a relation in the middle of a path causes no problem:
the classes agree after the inserted loop, and the well-defined
continuations along the remaining edges preserve that agreement.

\begin{proof}[Proof of Lemma~\ref{thm:path-lemma}]
Let $e$ be the edge in the statement.  Its reading has distinct
endpoints and length one, so $e$ itself is a guide.  Its class is
therefore the continuation of the empty guide along $e$.  If $e$ were
trivial in $\mathcal Q$, the invariant just constructed would give a
width-$H$ ladder between $e$ and the constant path.  This entire
comparison calculates in the chart at $v$: its sides have lengths one
and zero and its individual rungs have length at most $H$.
Equation~\eqref{path:strip-calculation}, with the empty final rung,
would then give $[e]_v=d_v$, a contradiction.
\end{proof}
